\documentclass[10pt]{article}
\usepackage{template}
\usepackage{amsmath,tikz}
\SetKeys{date=2026-11-24,vol=042,title={每日有趣新闻},author={ChatGPT 编辑部},note={短刊；观察期11月22—23日，数学拾趣含稍早预印本。}}
\begin{document}
\begin{documents}
\newsitem{COP31未能达成总体协议}{国际政治｜11月22—23日}{
联合国11月23日确认，安塔利亚气候大会未能形成最终气候协议。路透社报道，欧盟等参与方拒绝接受最终文本；化石燃料措辞成为主要分歧之一。联合国秘书长对结果表示失望。这是谈判结果，而非关于气候科学的新发现。\cite{un,reuters}

气候谈判面临一个典型的集体行动问题：减排收益跨越国界，而转型成本往往由本地企业、消费者和财政承担。各国的能源结构不同，因而同一项承诺可能对应完全不同的成本。对依赖油气出口的经济体而言，产业替代需要时间；对面临海平面上升的国家而言，延迟本身就是成本。气候融资可以缓和这种矛盾，但资金来源、偿还方式及使用监督同样需要协调。

卫报从气候行动角度报道，部分公正转型安排仍获同意。\cite{guardian} 因此，未能达成总体协议不等于所有既有义务或技术合作自动失效。下一步应观察正式决议清单、各国国内政策以及实际资金拨付，而不是仅凭会议结束时的政治评价判断减排趋势。
}
\newsitem{Roman望远镜开放第二轮研究申请}{航天｜11月23日NASA公告}{
NASA公布Roman空间望远镜第二轮研究申请安排，截止日期为2027年2月24日。此次申请面向广域仪器数据分析及相关理论、实验研究，而非新增观测任务；下一轮才预计开放新的通用天体物理巡天申请。\cite{roman}

开放数据与开放观测时间是两回事。前者让研究者能够复用既有观测，后者决定仪器实际指向何处以及积分时间如何分配。巡天数据的价值在于能够由不同团队使用不同统计模型重新分析，而不仅限于最初提出观测计划的团队。NASA另宣布计划于12月7日公布首批图像；在公布前不能将计划视作已经取得的科学成果。\cite{images}
}
\newsitem{数学拾趣：群的两种图结构}{组合与代数｜11月19日预印本，背景选读}{
设有限群为$G$，取逆封闭且不含单位元的生成集$S$。Cayley图以群元素为顶点，当$hg^{-1}\in S$时连接$g,h$。幂图则在一个元素属于另一个元素生成的循环子群时连边。两者分别表达乘法移动和循环子群中的包含关系。

Rabie与Iranmanesh研究连通Cayley图中同时属于幂图的边的最大比例，声称该比例为一当且仅当群为循环群，并给出若干有限群的具体结果。\cite{group} 本刊查阅了原文定义及摘要，未独立核验全部证明，同行评议状态未确认。

以四元Klein群为例，三个非单位元均为二阶，幂图是一颗三叶星；任意两个不同的非单位元构成生成集，对应Cayley图为四边形。四条边中只有两条属于幂图。若取全部三个非单位元，则Cayley图为完全图，六条边中三条属于幂图。因此最大比例为二分之一。这一例子由编辑部直接枚举，不依赖论文的普遍结论。

\begin{center}
\begin{tikzpicture}[scale=.8,every node/.style={circle,draw,inner sep=2pt}]
\node(e) at (0,0){$e$};\node(a) at (1.5,0){$a$};\node(b) at (0,1.5){$b$};\node(c) at (1.5,1.5){$ab$};
\draw[very thick,newsblue] (e)--(a) (e)--(b);
\draw[dashed] (a)--(c) (b)--(c);
\end{tikzpicture}

{\footnotesize 编辑部自绘，2026-11-24。Klein群的Cayley图；实线为同时属于幂图的边，虚线为其余边。}
\end{center}
}
\end{documents}
\bibliographystyle{plain}
\bibliography{refs}
\end{document}
